\paragraph{Effect of candidate-pool size.}
We extend the half-cycle diagnostic to $K\in\{4,8,16\}$ using nested
candidate pools and the same 2048 query--pool pairs per environment.
The $K=4$ pool consists of the first four candidates of each verified
$K=8$ pool, including the shared mismatched candidate. The $K=16$ pool
appends eight independent, uniformly sampled reference windows with
replacement (additional-candidate seed 67890). We retain the original
query windows, phase labels, and $K=8$ pools, and verify their OT costs
by recomputation. Thus the $K=8$ results are unchanged.
Coverage and rescue retain the same definitions with the selected
candidate taken from the corresponding $K$-candidate pool.

\begin{table}[H]
\centering\small
\begin{tabular}{lrccc}
\toprule
Environment & $K$ & Coverage (\%) & Rescue (\%) & Mean phase error \\
\midrule
hopper & 4 & 47.3 [44.9, 49.7] & 32.6 [30.2, 35.0] & 0.200 [0.190, 0.211] \\
hopper & 8 & 78.2 [76.5, 79.9] & 46.3 [43.6, 48.9] & 0.147 [0.139, 0.156] \\
hopper & 16 & 95.8 [94.9, 96.6] & 56.6 [54.1, 58.8] & 0.118 [0.112, 0.124] \\
ant & 4 & 51.8 [49.8, 53.9] & 40.1 [38.0, 42.3] & 0.168 [0.162, 0.174] \\
ant & 8 & 82.0 [80.2, 83.8] & 63.4 [60.9, 65.9] & 0.102 [0.098, 0.107] \\
ant & 16 & 97.7 [97.2, 98.3] & 81.4 [79.6, 83.3] & 0.073 [0.069, 0.078] \\
humanoid & 4 & 56.1 [54.4, 57.9] & 43.0 [41.3, 44.8] & 0.190 [0.184, 0.196] \\
humanoid & 8 & 85.2 [83.8, 86.7] & 64.0 [62.1, 65.9] & 0.121 [0.116, 0.125] \\
humanoid & 16 & 98.0 [97.4, 98.7] & 77.7 [76.0, 79.3] & 0.089 [0.085, 0.093] \\
snu\_humanoid & 4 & 55.8 [53.7, 58.0] & 48.0 [45.9, 50.4] & 0.168 [0.163, 0.174] \\
snu\_humanoid & 8 & 85.1 [83.8, 86.5] & 72.0 [70.1, 73.9] & 0.095 [0.092, 0.099] \\
snu\_humanoid & 16 & 97.9 [97.4, 98.5] & 87.2 [85.6, 88.8] & 0.058 [0.055, 0.062] \\
\bottomrule
\end{tabular}
\caption{Nested-pool phase-rescue comparison at a target half-cycle mismatch. Entries are point estimates [95\% trajectory-cluster bootstrap CI], using 512 held-out query windows and four pools per window. $K=4$ uses the first four candidates of the verified $K=8$ pool; $K=16$ appends eight uniformly sampled candidates. All conditions share the same mismatched baseline.}
\label{tab:phase_rescue_k_sweep}
\end{table}


Increasing $K$ improves coverage and rescue rates and reduces mean
selected phase error in all four environments. From $K=8$ to $K=16$,
rescue increases by 10.3, 18.0, 13.7, and 15.1 percentage points for
Hopper, Ant, Humanoid, and SNU Humanoid, respectively. Paired 95\%
trajectory-cluster bootstrap intervals for these increases are
$[8.5,12.3]$, $[16.4,19.7]$, $[12.0,15.4]$, and $[13.2,17.1]$ percentage
points. These intervals are computed on within-pair differences,
using the same trajectory-cluster resampling procedure, and are
pointwise without multiplicity correction. The phase-error reduction
from $K=8$ to $K=16$ is smaller than from $K=4$ to $K=8$ in every
environment. These results concern minimum-cost selection on held-out
expert windows, not policy-training performance.
